% Independent derivation; no production implementation is imported.
\subsection{Exact two-prefix adaptive width and a batch precision certificate}
Let $y\in\mathbb R_+^n$ satisfy $\sum_{i=1}^n y_i=E$, where $E\geq0$ is
known, and let the audited contrast be $T(y)=c^\top y+d$.  Each available query
returns an exact prefix total $P_k(y)=\sum_{i=1}^k y_i$, $1\leq k<n$.
Queries have unit cost and every execution may make at most two queries.
For indices $a\leq b$, write
\[
 r(a,b)=\max_{a\leq i\leq b}c_i-\min_{a\leq i\leq b}c_i.
\]
The ranges of singleton segments are zero.  If queried prefixes partition
$\{1,\ldots,n\}$ into consecutive segments $I$, their masses $m_I$ are known
and the sharp interval width is
\begin{equation}\label{eq:counter_segment_width}
 W_Q(y)=\sum_I m_I r(I).
\end{equation}
Indeed, the nonnegative targets within a segment can place all its mass on a
minimum or maximum coefficient, independently across segments.  These
completions attain both endpoints.

After first querying prefix $k$ and observing $t\in[0,E]$, the exact worst
terminal width if the second query is $j\ne k$ is the affine function
\begin{equation}\label{eq:counter_conditional_line}
 g_{kj}(t)=
 \begin{cases}
 t\max\{r(1,j),r(j+1,k)\}+(E-t)r(k+1,n),&j<k,\\
 tr(1,k)+(E-t)\max\{r(k+1,j),r(j+1,n)\},&j>k.
 \end{cases}
\end{equation}
For example, when $j<k$, the first segment's known mass $t$ can be split
arbitrarily between the new segments.  Maximizing its contribution therefore
places the mass in the segment with larger coefficient range.  The unsplit
right segment contributes $(E-t)r(k+1,n)$.  This proves equality, rather than
only an upper bound; the case $j>k$ is symmetric.

Consequently, the optimal deterministic adaptive two-query worst-width value
is
\begin{equation}\label{eq:counter_adaptive_value}
 A=\min_{1\leq k<n}\max_{0\leq t\leq E}\min_{j\ne k}g_{kj}(t).
\end{equation}
For a fixed first cut, the policy chooses a minimizing second cut after
observing $t$.  Conversely, any deterministic policy must choose such a
second cut, and the conditional maximum in
\eqref{eq:counter_conditional_line} is attainable.  Maximizing over $t$ and
then minimizing over the first cut proves \eqref{eq:counter_adaptive_value}.
Stopping after one query cannot improve width over a legal additional query,
so allowing early stopping does not reduce this value.  The lower envelope
of finitely many affine functions is concave and piecewise affine.  Its
maximum is attained at an endpoint or at a pairwise line intersection in
$[0,E]$.  Equivalently, it is the optimum of
$\max\{\eta:0\leq t\leq E,\ \eta\leq g_{kj}(t)\ \forall j\ne k\}$.
For rational input, enumerating these candidates gives an exact rational
value.  This procedure optimizes a hard-budget worst-width guarantee; it
does not optimize expected requests until certification.

For a prescribed threshold $\tau\geq0$, take any finite collection of
feasible profiles $y^1,\ldots,y^s$ and enumerate all two-prefix batches $Q$.
Define the exact binary matrix
$H_{\ell Q}=\mathbf1\{W_Q(y^\ell)\leq\tau\}$.
The finite game
\begin{equation}\label{eq:counter_batch_upper}
 U(\tau)=\max_{\pi\in\Delta(\mathcal Q)}\min_\ell
             \sum_Q\pi_QH_{\ell Q}
 =\min_{w\in\Delta_s}\max_Q\sum_\ell w_\ell H_{\ell Q}
\end{equation}
is an upper bound on the worst-record success probability achievable by
\emph{any randomized nonadaptive batch}.  Restricting the adverse record
set can only increase that success value.  For any feasible adverse weights
$w$, the explicitly calculated quantity
$\max_Q\sum_\ell w_\ell H_{\ell Q}$ is already a valid upper bound, whether
or not $w$ is optimal.  Smaller batches are dominated by their extensions
to two queries.  In general, a concentration-profile game is not an exact
solution of the continuum precision problem: threshold success need not
attain its worst value at a vertex.  No such claim is needed for the upper
bound.  If $\tau\geq A$, the adaptive procedure succeeds for every feasible
record, and $1-U(\tau)$ is a certified lower bound on its advantage in
worst-record precision success over every randomized batch.

\paragraph{An exact six-hour separation.}
Let $E=1$, $c=(0,0,1,1,0,1)$ and $\tau=1/2$.  Query prefix 4 first and write
$t=P_4(y)$.  Query prefix 5 if $t\leq1/2$, and prefix 2 otherwise.  Its
terminal width is $\min(t,1-t)\leq1/2$.  Formula
\eqref{eq:counter_adaptive_value} gives first-cut values
$(1,1,1,1/2,1)$, so its optimal deterministic adaptive worst width is exactly
$A=1/2$.

The fixed batch $(2,4)$ has width $1-t$; the fixed batch $(4,5)$ has width
$t$.  Their equal mixture succeeds with probability at least $1/2$ for
every record.  Conversely, put adverse probability $1/2$ on $e_1$ and
$1/2$ on $e_5$.  At $e_5$, every batch except $(4,5)$ has width 1, since
isolating the fifth coefficient from its two neighbors requires both cuts
4 and 5.  The batch $(4,5)$ has width 1 at $e_1$.  Thus every batch succeeds
on at most one of these two adverse records, and every randomized batch
has worst-record success at most $1/2$.  Its optimal success is therefore
exactly $1/2$, whereas the adaptive success is 1.

These policies nevertheless have identical optimal worst expected width,
$1/2$.  The equal batch mixture has expected width $1/2$ at every record;
the common record $y^*=(e_1+e_5)/2$ has width at least $1/2$ for every final
batch, including a batch selected by an adaptive transcript.  Exact
prefix answers and a fixed internal seed define the same observation fiber
as the final set of queried prefixes: every completion agreeing with those
answers follows the same adaptive query and stopping sequence.  Thus the
common-record width lower bound applies to adaptive policies as well.
The precision advantage concerns reliability of the terminal interval's
width, not winner identification, forecasting error, or retrieval cost.

\paragraph{Comparator discipline.}
A natural one-step minimax greedy rule first minimizes
$E\max\{r(1,k),r(k+1,n)\}$ and then, given $t$, minimizes
$g_{kj}(t)$.  All five first cuts tie under this greedy first-stage
criterion in the example.  Choosing cut 1 by index gives worst width 1,
while choosing cut 4 gives $1/2$.  Consequently, an apparent factor-two
gain over lexicographic greedy would be a tie-selection effect.  The
separation proved above instead uses every randomized batch as comparator.
An empirical greedy study should freeze its tie rule, give all methods
identical available prefixes and totals, and report sensitivity to the
best first-stage greedy tie when possible.
